\documentclass[11pt]{article}

% ---- packages ----
\usepackage[letterpaper,margin=1in]{geometry}
\usepackage{microtype}
\usepackage{amsmath}
\usepackage{booktabs}
\usepackage{array}
\usepackage{graphicx}
\usepackage{xcolor}
\usepackage[most]{tcolorbox}
\usepackage{enumitem}
\usepackage{caption}
\usepackage[hidelinks]{hyperref}
\usepackage{parskip}
\usepackage{tocloft}
\setlength{\cftbeforesecskip}{0pt}
\usepackage[nameinlink]{cleveref}

\graphicspath{{figures/}}
\captionsetup{font=small,labelfont=bf}
\setlist{leftmargin=1.4em,topsep=3pt,itemsep=2pt}

% ---- teaching / takeaway callout boxes ----
\definecolor{bgblue}{HTML}{2563EB}
\definecolor{bggreen}{HTML}{15803D}
\newtcolorbox{background}[1]{
  breakable, colback=bgblue!5, colframe=bgblue!60, coltitle=bgblue!75!black,
  fonttitle=\bfseries\small, title={Background you'll need --- #1},
  boxrule=0.6pt, left=6pt, right=6pt, top=4pt, bottom=4pt, arc=2pt}
\newtcolorbox{takeaway}[1]{
  breakable, colback=bggreen!6, colframe=bggreen!55, coltitle=bggreen!60!black,
  fonttitle=\bfseries\small, title={#1},
  boxrule=0.6pt, left=6pt, right=6pt, top=4pt, bottom=4pt, arc=2pt}

\newcommand{\code}[1]{\texttt{\small #1}}

\title{\textbf{The Benchmark Chooses the Winner}\\[2pt]
\large An Annotated, Plain-Language Edition\\[2pt]
\normalsize\itshape Measuring what fine-tuning does to small safety guards}
\author{Reza Rahimi, PhD \\ \small JazzX AI, Los Altos, CA \\ \small\texttt{reza.rahimi@jazzx.ai}}
\date{}

\begin{document}
\maketitle

\begin{abstract}
\noindent
Teams often take a small chat model, do a quick fine-tune to turn it into a
\emph{safety guard} (a filter that flags harmful prompts), and report how well it scores
on a public benchmark. This edition re-tells a study that asks a sharper question:
\emph{what did the fine-tune change compared to the exact same model before fine-tuning,
and does that change hold up on datasets the model never trained on?} On four small
instruction models (1.5--4B parameters), each fine-tuned five times with LoRA, we observe a
split. On test data from the \textbf{same datasets used in training} (``represented''
sources), the observed fine-tuning gain is large: the main score rises by about
\textbf{$+0.32$} (to $\approx0.98$ for every model). On \textbf{four dataset-held-out
transfer benchmarks} not used in training, the observed average change is a small loss
(about \textbf{$-0.06$}), and at a calibration-selected diagnostic threshold it catches
18 percentage points fewer harmful examples on a hard
held-out harm set (HarmBench recall $78.0\%\!\to\!60.0\%$). Crucially, the average hides a
checkpoint split: SmolLM2 improves, Qwen2.5's interval crosses zero, and SmolLM3/Qwen3 decline.
With four same-row checkpoint contrasts, this does not establish base strength as a predictor. The
lesson --- and the title --- is that \textbf{whichever benchmark you choose to report decides
whether fine-tuning ``looks like'' a win}. A clearly separated section then previews a
\emph{preliminary clean-v2 retrospective Paper B analysis}: averaging the base and tuned guard's temperature-scaled
scores (a two-model ``ensemble'') recovers some transfer ranking in this fixed panel. This is
an exploratory follow-up, not a Paper A finding or an established remedy. Every concept
needed to read this (guard scoring, LoRA, average precision, the bootstrap, calibration,
ensembling) is taught inline. The Paper A numbers are retrospective fixed-panel estimates with
conditional resampling intervals, not prospective population claims. Re-running locked code
cannot make already-inspected transfer sets prospective; that would require a new sealed
cohort. This is a plain-language companion to the formal paper, which remains the source
of truth wherever this edition simplifies.
\end{abstract}

\setcounter{tocdepth}{1}
\begingroup
\small
\tableofcontents
\endgroup

\clearpage
% ============================================================
\section{Introduction}
% ============================================================

A \textbf{prompt-safety guard} reads a user's prompt and decides whether it is okay
(\textbf{safe}) or an attack or harmful request (\textbf{unsafe}). ``Unsafe'' here means one
of three things: \emph{harmful content} (asking for something dangerous), a \emph{jailbreak}
(trying to trick the model into ignoring its rules), or a \emph{prompt injection} (hiding
instructions in the input to hijack the system). Guards are attractive because they are
small and cheap to run in front of a bigger model.

The common recipe is: take a general chat model, fine-tune it on your own labeled prompts,
and report its score on a benchmark. The trouble is that a benchmark score, on its own,
hides the quantity a practitioner actually needs: \emph{what changed relative to the same
model before fine-tuning, and does that change generalize?} If you only compare your tuned
guard against \emph{other} models, a gain on the datasets you trained on can masquerade as a
broadly smarter guard.

\begin{background}{what a guard is and how we score it}
\small
The guard is just a language model asked to answer with one word: \code{safe} or
\code{unsafe}. Instead of letting it write a sentence, we do \textbf{one forward pass} and
look only at the two candidate next-words. For each, the model produces a raw preference
number (a \textbf{logit}). The score is the gap between them,
\[
s(x) = z_{\text{unsafe}} - z_{\text{safe}},
\]
positive meaning ``leans unsafe.'' A \textbf{softmax} over just those two turns the pair into
bounded shares that add up to 1; the unsafe share is a score, not automatically a calibrated
probability. \emph{Mini-example:} if $z_{\text{unsafe}}=2.0$ and
$z_{\text{safe}}=1.0$, then $s=1.0$ and
$e^{2.0}/(e^{2.0}+e^{1.0}) \approx 0.73$. This is an unsafe score of 0.73, not necessarily a
literal 73\% chance. Ranking prompts by this score is all AP needs.
\end{background}

\paragraph{The question.} For a fixed panel of four checkpoints and one frozen LoRA recipe,
does fine-tuning change a guard's discrimination (how well it separates unsafe prompts from
safe ones) \emph{differently} on datasets represented during training versus datasets held out
from training?

\paragraph{Contributions.} (1) A controlled, same-model comparison: every fine-tuned guard is
measured against \emph{its own} untuned base on identical prompts, metrics, and scoring.
(2) A separation of two effects that are usually blended into one leaderboard number: gains
on \emph{represented} sources versus \emph{transfer} to held-out datasets, with per-model and
per-benchmark detail. (3) A fully auditable release: hashes, seed-level results, the scoring
table, and generated tables/figures.

\paragraph{Scope.} English input-prompt classification only; a binary safe/unsafe decision;
small models (1.5--4B); one LoRA recipe; three represented datasets; four transfer datasets;
two single-class stress sets; four fixed checkpoints. It is \emph{not} a leaderboard, a
fine-tuning-objective comparison, an ensembling or fairness study, or a domain-compliance
case study. The later composition section is explicitly a separate preliminary Paper B preview.

% ============================================================
\section{Background you'll need}
% ============================================================
This section teaches the five ideas the results rest on. If you already know average
precision, the bootstrap, and calibration, skip to Section~\ref{sec:design}.

\subsection{LoRA supervised fine-tuning (SFT)}
\begin{background}{LoRA-SFT}
\small
LoRA fine-tuning~\cite{hu2021lora} does \textbf{not} create a new model. The original model is \textbf{frozen}
(no weight changes). Instead you bolt on a small set of extra trainable weights --- an
\textbf{adapter} --- plugged into specific matrices inside the model. Only the adapter learns,
so it is cheap: you train a small fraction of extra parameters (millions, not billions) and
ship a small file that snaps onto the frozen base. ``Supervised'' means training on labeled
examples (prompt $+$ correct verdict word). \textbf{Completion-only loss on the verdict
token} means the model is graded \emph{only} on producing the single answer token, not on
echoing the prompt --- so all the learning pressure lands on that one decision.
\end{background}

\subsection{Average precision (AP), and why not accuracy}
\label{subsec:ap}
\begin{background}{average precision}
\small
The guard outputs a \emph{score} (the logit difference), and a good guard gives
\emph{higher} scores to truly-unsafe prompts. Two everyday quantities describe how it does:
\textbf{precision} $=$ of the prompts flagged so far, the fraction that are truly unsafe; and
\textbf{recall} $=$ of all the unsafe prompts, the fraction we have flagged.
\textbf{Average Precision (AP)} rolls these up with \textbf{no threshold to pick}: it is the
area under the precision--recall curve (0--1, higher better) --- in words, how well the guard
\emph{ranks} unsafe prompts above safe ones. Why not plain \textbf{accuracy}? Accuracy needs a
fixed cutoff and can be a trap --- if 95\% of prompts are safe, a guard that labels
\emph{everything} safe scores 95\% accuracy while catching zero attacks. AP does not fall for
that. \emph{Mini-example:} two unsafe and three safe prompts, ranked
\{unsafe, safe, unsafe, safe, safe\}. Precision at the first unsafe $=1/1=1.0$ (1 unsafe among
the top 1); at the second $=2/3\approx0.67$ (2 unsafe among the top 3);
$\text{AP}=(1.0+2/3)/2\approx\mathbf{0.83}$ (a perfect ranking scores 1.0).
\emph{Tie-aware} means two prompts with the same score get fair, not lucky, credit.
\end{background}
\noindent Throughout, the headline metric is \textbf{benchmark-macro AP}: average the AP
across the datasets in a group so each dataset counts equally (``macro''), then average across
the four models, and for fine-tuned guards across the five seeds too.

\subsection{Represented sources vs.\ held-out transfer}
\label{subsec:regimes}
\begin{background}{the represented-vs-transfer distinction (the crux)}
\small
\textbf{Represented} sources are fresh, held-out rows from the \emph{same} datasets used in
training (same style and labeling habits, different examples). Like taking this year's exam
on a topic you crammed from last year's --- you do well, partly because you learned the test's
patterns. \textbf{Transfer} is a \emph{new-material exam}: four entirely different datasets
never used in training. Reporting a single mixed number hides the difference between ``got
better at \emph{these} datasets'' and ``got broadly better'' --- exactly the confusion this
paper is about. \emph{(Honest caveat: the transfer sets were looked at during development, so
they are ``dataset-held-out,'' not a sealed surprise.)}
\end{background}

\subsection{Conditional fixed-panel bootstrap intervals}
\label{subsec:boot}
\begin{background}{bootstrap confidence intervals}
\small
Each headline number is one estimate from the fixed models, datasets, and seeds we ran. To
describe how it changes under a specified resampling scheme, we \textbf{resample our own
results} (with replacement) 10{,}000 times and recompute the average each time (a
\textbf{paired hierarchical bootstrap}: \emph{paired} = base vs.\ fine-tuned, to measure the
\emph{change}; \emph{hierarchical} = resample at two levels, which seeds count and --- because
the eval sets contain many near-duplicate prompts --- how much each group of duplicates
counts; the four models are held fixed). The middle 95\% of those recomputed averages is the
\textbf{two-sided 95\% percentile-bootstrap interval}; one-sided LCB/UCB values are the
corresponding lower/upper bootstrap quantiles. These are conditional fixed-panel summaries:
the four checkpoints, benchmark collection, labels, and analysis choices remain fixed. They
are not 95\% probability statements about a population effect or guarantees for future models
and data. The study is therefore \textbf{descriptive}: it reports resampling precision but no
significance test or pass/fail verdict. A locked rerun can strengthen execution provenance;
prospective evidence requires a new sealed cohort and a prospectively locked analysis.
\end{background}

\subsection{Calibration and the operating point}
\label{subsec:calib}
\begin{background}{calibration, thresholds, TPR and FPR}
\small
AP is threshold-free, but a yes/no diagnostic needs a cutoff. First,
\textbf{calibration (temperature scaling)} divides the logit score by one number fitted on the
held-back calibration set. This does not reorder prompts and may improve frequency alignment
on that distribution; it does not guarantee calibrated probabilities on transfer data. Second,
the \textbf{threshold} maximizes calibration recall among cutoffs whose one-sided 95\%
row-level Clopper--Pearson upper bound on pooled calibration-negative FPR is at most 5\%. This is a
calibration-sample diagnostic constraint, not a production or distribution-shift guarantee.
Then the frozen cutoff is \emph{measured} on test data: \textbf{TPR} (recall) $=$ share of
unsafe prompts caught (higher better); \textbf{FPR} $=$ share of safe prompts wrongly flagged
(lower better). \emph{Illustrative mini-example (not the study's threshold calculation):}
4 false alarms among 80 safe prompts gives an empirical FPR of 5\%; catching 15 of 20 unsafe
prompts gives recall of 75\%. The study additionally requires the one-sided FPR upper bound,
not just the empirical rate, to meet the calibration criterion.
\end{background}

% ============================================================
\section{Study design}
\label{sec:design}
% ============================================================

\subsection{The panel and the frozen recipe}
Four instruction models are each adapted with five training seeds (42--46), giving 20
fine-tuned guards; the four untuned bases are scored once and reused ($4+20=24$ ``bundles''
in total, and $79{,}392$ scored rows). Every cell uses the identical recipe in
Table~\ref{tab:panel}. The training seed controls adapter initialization and dropout, while a
separate fixed data-order seed keeps row order identical across runs. Five seeds describe
run-to-run optimization variation under this one recipe; they do not measure variation across
new checkpoints, datasets, or recipes.

\begin{table}[t]
\centering\small
\caption{The fixed four-checkpoint panel and the one LoRA-SFT recipe used for every cell.}
\label{tab:panel}
\begin{tabular}{lll}
\toprule
Checkpoint & Params & Role \\
\midrule
Qwen2.5-1.5B-Instruct & 1.5B & base $\to$ 5 SFT seeds \\
SmolLM2-1.7B-Instruct & 1.7B & base $\to$ 5 SFT seeds \\
SmolLM3-3B~\cite{bakouch2025smollm3}            & 3B   & base $\to$ 5 SFT seeds \\
Qwen3-4B             & 4B   & base $\to$ 5 SFT seeds \\
\midrule
\multicolumn{3}{l}{\emph{Recipe:} LoRA $r{=}32$, $\alpha{=}64$, dropout $0.05$ on
q,k,v,o,gate,up,down; completion-only loss;}\\
\multicolumn{3}{l}{\phantom{\emph{Recipe:}} $1{,}200$ training rows; 300 steps; learning rate
$2\times10^{-4}$ cosine; effective batch 4; max length 1024.}\\
\bottomrule
\end{tabular}
\end{table}

\subsection{Training data, label mapping, and decontamination}
Training reads one frozen manifest of \textbf{1{,}200 rows}: 400 from each of three datasets
--- ToxicChat~\cite{lin2023toxicchat}, Prompt-Injections, and Jailbreak-Classification --- with 200 safe and 200
unsafe per dataset, selected by a fixed hash rank so the row set is identical across every
model and seed. Each dataset's native label is mapped to safe/unsafe by a fixed rule (e.g.\
ToxicChat \code{toxicity=1}$\to$unsafe; Prompt-Injections \code{label=1}$\to$unsafe;
Jailbreak-Classification \code{type} starting ``jailbreak''$\to$unsafe). These are
\emph{different} native policies, not one universal definition of ``unsafe.''

To reduce train/evaluation contamination, the builder removes exact matches and MinHash
near-duplicate families at the prespecified threshold, refusing to let a detected family
straddle the train/evaluation boundary. This is an auditable heuristic screen, not proof that
no semantic overlap remains.

\begin{table}[t]
\centering\small
\caption{Datasets and their roles. Represented sources appear (as different rows) in both
training and the represented test; transfer and stress sets are test-only.}
\label{tab:data}
\begin{tabular}{lll}
\toprule
Group & Datasets & Use \\
\midrule
Represented & ToxicChat, Prompt-Injections, Jailbreak-Classification & train $+$ represented test \\
Transfer & JailbreakBench, XSTest, WildGuardTest, WildJailbreak & test only (never trained on) \\
Stress & OR-Bench (all benign), HarmBench (all harmful) & one-sided diagnostics only \\
\bottomrule
\end{tabular}
\end{table}

\subsection{The three test groups}
\begin{itemize}
\item \textbf{Represented} $=$ held-out rows from the three training datasets --- measures
discrimination on \emph{sources seen in training}.
\item \textbf{Transfer} $=$ four datasets withheld from training~\cite{chao2024jailbreakbench,rottger2024xstest,han2024wildguard,jiang2024wildteaming}
--- measures whether the skill \emph{generalizes} (Section~\ref{subsec:regimes}).
\item \textbf{Stress} $=$ two single-class sets. OR-Bench~\cite{cui2024orbench} is all benign (measures false
alarms); HarmBench~\cite{mazeika2024harmbench} is all harmful (measures catch rate). Because each has only one class,
\textbf{no AP is computed on them} --- they are diagnostics.
\end{itemize}

\subsection{What we compute}
For a group $R$, checkpoint $b$, and seed $r$, the metric is the benchmark-macro AP
(Section~\ref{subsec:ap}). The change from fine-tuning is the paired difference
\[
\Delta_R(b,r) = \text{macroAP}_R(\text{SFT}_{b,r}) - \text{macroAP}_R(\text{base}_b),
\]
and the panel aggregate averages the four checkpoints (averaging seeds within each):
\[
\bar\Delta_R = \frac{1}{4}\sum_b\Big[\tfrac{1}{5}\sum_r \text{macroAP}_R(\text{SFT}_{b,r}) -
\text{macroAP}_R(\text{base}_b)\Big].
\]
The primary result is the \emph{pair}
$(\bar\Delta_{\text{represented}}, \bar\Delta_{\text{transfer}})$ --- we keep both numbers
rather than collapsing them into one ``specialization score.'' Uncertainty is the paired
hierarchical bootstrap of Section~\ref{subsec:boot}; the analysis mode is
\emph{descriptive} (no significance test).

% ============================================================
\section{Results}
% ============================================================

\subsection{Represented sources: a large observed gain}
On the datasets the models trained on, every fine-tuned guard reaches ${\approx}0.98$
macro-AP regardless of where its base started (Table~\ref{tab:primary}, left). The panel
aggregate change is $\bar\Delta_{\text{represented}} = \mathbf{+0.3234}$ (95\% CI
$[+0.2647, +0.3690]$). Notice that the \emph{weakest} base gains the most (SmolLM2:
$0.4524\!\to\!0.9806$, $+0.5282$) and the \emph{strongest} gains the least (Qwen3-4B:
$0.8855\!\to\!0.9837$, $+0.0981$) --- largely a ceiling effect: everyone ends near 0.98, so
whoever started lowest had the most room to climb. Per-dataset gains: ToxicChat $+0.1880$,
Prompt-Injections $+0.3704$, Jailbreak-Classification $+0.4119$ (Table~\ref{tab:perbench}).

\subsection{Transfer: a small average loss that hides a split}
On the four never-trained-on datasets, the same fine-tuning is a small average loss:
$\bar\Delta_{\text{transfer}} = \mathbf{-0.0589}$ (95\% CI $[-0.0837, -0.0321]$). But the
per-checkpoint column of Table~\ref{tab:primary} shows heterogeneous descriptive estimates:
one is positive, one interval crosses zero, and two are negative. The losses concentrate on the
jailbreak-style sets (WildJailbreak $-0.0792$, JailbreakBench $-0.0776$), while the
over-refusal set barely moves (XSTest $-0.0120$); see Table~\ref{tab:perbench}.

\begin{table}[t]
\centering\small
\caption{Primary result: base and mean-SFT benchmark-macro AP per checkpoint, with the paired
change $\Delta$ and its two-sided 95\% bootstrap interval. Positive $\Delta$ $=$ fine-tuning
helped. The last row is the panel aggregate.}
\label{tab:primary}
\setlength{\tabcolsep}{3pt}
\begin{tabular}{l rr r@{\ }l rr r@{\ }l}
\toprule
& \multicolumn{3}{c}{Represented} && \multicolumn{3}{c}{Transfer} &\\
\cmidrule(lr){2-4}\cmidrule(lr){6-8}
Checkpoint & base & SFT & $\Delta$ & [95\% CI] & base & SFT & $\Delta$ & [95\% CI]\\
\midrule
Qwen2.5-1.5B & 0.6334 & 0.9878 & $+0.3544$ & [0.2731, 0.4150] & 0.8187 & 0.7798 & $-0.0389$ & [$-0.0829$, $+0.0062$]\\
SmolLM2-1.7B & 0.4524 & 0.9806 & $+0.5282$ & [0.4555, 0.5748] & 0.7904 & 0.8304 & $+0.0400$ & [0.0003, 0.0776]\\
SmolLM3-3B  & 0.6621 & 0.9751 & $+0.3130$ & [0.2419, 0.3701] & 0.9102 & 0.8234 & $-0.0869$ & [$-0.1114$, $-0.0613$]\\
Qwen3-4B   & 0.8855 & 0.9837 & $+0.0981$ & [0.0545, 0.1479] & 0.9438 & 0.7939 & $-0.1499$ & [$-0.1963$, $-0.1050$]\\
\midrule
\textbf{Panel aggregate} & -- & -- & $+0.3234$ & [0.2647, 0.3690] & -- & -- & $-0.0589$ & [$-0.0837$, $-0.0321$]\\
\bottomrule
\end{tabular}
\end{table}

\begin{table}[t]
\centering\small
\caption{Per-dataset paired AP change (base $\to$ SFT). Represented gains are large and
uniform; transfer losses concentrate on jailbreak-style sets.}
\label{tab:perbench}
\begin{tabular}{llr}
\toprule
Group & Dataset & $\Delta$ AP \\
\midrule
Represented & ToxicChat & $+0.1880$ \\
Represented & Prompt-Injections & $+0.3704$ \\
Represented & Jailbreak-Classification & $+0.4119$ \\
\midrule
Transfer & JailbreakBench & $-0.0776$ \\
Transfer & XSTest & $-0.0120$ \\
Transfer & WildGuardTest & $-0.0669$ \\
Transfer & WildJailbreak & $-0.0792$ \\
\bottomrule
\end{tabular}
\end{table}

\subsection{The key nuance: checkpoint heterogeneity behind the average}
Read the ``SFT'' columns of Table~\ref{tab:primary}: on transfer, fine-tuning squeezes every
model into a narrow $\approx0.78$--$0.83$ band, even though the bases ranged from 0.79 to 0.94.
The four point estimates therefore have different directions. But $\Delta=\mathrm{SFT}-\mathrm{base}$,
so relating that change to base AP measured on the same rows is mathematically coupled. Four
checkpoints cannot establish a fixed destination or a base-competence law.

\begin{takeaway}{Why the $-0.06$ average is misleading}
\small
As a toy example, imagine two changes of $+0.10$ and $-0.15$. Their average is only $-0.025$,
which hides both opposing effects. In the observed data, SmolLM2 gains ($+0.0400$), Qwen2.5's
interval crosses zero, and SmolLM3 ($-0.0869$) and Qwen3-4B ($-0.1499$) lose. These are
clean-v2 retrospective checkpoint estimates, not a validated predictor based on starting strength.
\end{takeaway}

\subsection{The specialization plane}
Figure~\ref{fig:plane} puts both changes on one chart: each dot is one (checkpoint, seed);
right means represented AP went up, up means transfer AP went up. \textbf{15 of the 20 dots}
land in the lower-right ``specialization'' quadrant (better on familiar data, worse on new);
the other 5 are in ``uniform gain'' (SmolLM2 seeds 42--45 plus Qwen2.5 seed 42), and none
are in a loss quadrant. The dominant story is specialization.

\begin{figure}[t]
\centering
\includegraphics[width=0.82\linewidth]{specialization_plane.pdf}
\caption{The specialization plane. Horizontal $=$ change on represented (trained-on) sources;
vertical $=$ change on transfer (new) datasets. Each dot is one (checkpoint, seed); the
$\times$ is the panel mean. All 20 dots move right (represented up); 15 also move down
(transfer down).}
\label{fig:plane}
\end{figure}

\subsection{At a calibration-selected diagnostic threshold}
AP is threshold-free; Table~\ref{tab:op} shows what happens at the calibration-selected
operating point from Section~\ref{subsec:calib}. On represented data the
guard becomes far more useful --- recall jumps $13.0\%\!\to\!76.9\%$ at only ${\sim}1\%$ false
alarms. On transfer data recall rises ($51.7\%\!\to\!58.1\%$) while false alarms get
\emph{worse}: benchmark-macro FPR $8.1\%\!\to\!15.5\%$, pooled FPR $4.3\%\!\to\!17.0\%$
(\emph{pooled} $=$ every test row combined into one set, so larger datasets dominate;
\emph{macro} $=$ each dataset weighted equally --- which is why the base's pooled $4.3\%$ and
macro $8.1\%$ differ).

\begin{takeaway}{Why transfer FPR (15.5\%) can exceed the 5\% calibration criterion}
\small
The 5\% criterion constrains a row-level one-sided bound computed from pooled \emph{calibration}
negatives. It is not a guarantee for a shifted test distribution. The cutoff is frozen and
never re-tuned, so realized transfer FPR can exceed 5\% (the base is already at 8.1\%).
\end{takeaway}

\subsection{Stress sets: false alarms and missed attacks}
The two single-class diagnostics (Table~\ref{tab:op}, lower panel) sharpen the picture.
Fine-tuning leaves over-blocking of benign prompts essentially flat (OR-Bench FPR
$11.8\%\!\to\!12.0\%$, a 0.2-point increase). But on hard held-out harmful prompts it catches
\emph{fewer}: HarmBench recall falls $78.0\%\!\to\!60.0\%$, an observed
\textbf{18.0-point drop}. That is the largest observed downside in the study and is
safety-relevant: the specialized guard misses about 18 percentage points more harmful examples
on this stress set. A one-class diagnostic does not estimate performance on mixed deployment
traffic.

\begin{table}[t]
\centering\small
\caption{Calibration-selected threshold behavior (row-level one-sided 95\% calibration FPR upper bound
at most 5\%) and single-class stress diagnostics. TPR $=$ recall; FPR $=$ false-alarm rate;
$N$ $=$ rows in the relevant class. This calibration criterion is not a transfer guarantee.}
\label{tab:op}
\begin{tabular}{ll rr r r}
\toprule
Group & Measure & base & SFT & change & $N$ \\
\midrule
Represented & TPR (recall) & 13.0\% & 76.9\% & $+63.9$ & 313 \\
Represented & FPR (macro)  & 1.7\%  & 1.1\%  & $-0.6$  & 364 \\
Transfer    & TPR (recall) & 51.7\% & 58.1\% & $+6.4$  & 790 \\
Transfer    & FPR (macro)  & 8.1\%  & 15.5\% & $+7.4$  & 790 \\
Transfer    & FPR (pooled) & 4.3\%  & 17.0\% & $+12.7$ & 790 \\
\midrule
Stress: OR-Bench & benign FPR & 11.8\% & 12.0\% & $+0.2$ & 400 \\
Stress: HarmBench & recall    & 78.0\% & 60.0\% & $-18.0$ & 200 \\
\bottomrule
\end{tabular}
\end{table}

% ============================================================
\section{Separate Paper B analysis: compose, don't tune}
\label{sec:fix}
% ============================================================
This section is \textbf{not a Paper A result}. It previews a separate, clean-v2 retrospective Paper B
analysis, ``Compose, Don't Tune,'' that reuses the underlying Paper A scored rows to explore a
candidate mitigation. It asks whether keeping the original untuned model in the decision may
recover some transfer ranking after fine-tuning.

\begin{background}{combining two guards (an ``ensemble'')}
\small
Run \emph{two} guards on each prompt --- the untuned base and the fine-tuned adapter --- and
average their separately temperature-scaled unsafe scores. (Using more than one model and combining
their outputs is called an \textbf{ensemble}.) Think of it as a \emph{second opinion}: the
two guards can preserve different ranking signal when their errors differ. The guards share the
same backbone with the adapter switched on or off, so composition needs \textbf{no extra training},
but it still doubles the forward-pass work per prompt unless the passes run in parallel.
\emph{Mini-example:} on a new prompt the base scores $0.30$
unsafe and the tuned guard scores $0.90$; the ensemble scores $(0.30+0.90)/2 = 0.60$.
\end{background}

In this preliminary Paper B analysis (\Cref{tab:compose}), the ensemble keeps almost all of
the tuned guard's observed gain on familiar data ($0.982 \to 0.962$), while its transfer point
estimate is higher than the tuned guard's ($0.807 \to 0.883$), pulling the score toward the
untuned base without extra fine-tuning. On transfer, composed minus SFT is $+0.075$ (95\% CI
$[+0.058,+0.093]$) and composed minus base is $+0.017$ (95\% CI $[+0.005,+0.030]$).
On represented data, composition is $-0.019$ below SFT (95\% CI $[-0.031,-0.010]$), so this
is not a universal or Pareto improvement.

\begin{table}[t]
\centering\small
\caption{Separate preliminary Paper B composition analysis. Benchmark-macro AP (0--1, higher better),
averaged over the four checkpoints. The ensemble keeps the tuned guard's familiar-data gain and
recovers most of its lost new-data score.}
\label{tab:compose}
\begin{tabular}{lrr}
\toprule
Guard & Familiar (represented) & New (transfer) \\
\midrule
Untuned base & 0.658 & 0.866 \\
Tuned (SFT) & \textbf{0.982} & 0.807 \\
Composed (base $+$ tuned, averaged) & 0.962 & \textbf{0.883} \\
\bottomrule
\end{tabular}
\end{table}

\begin{takeaway}{What the preliminary point estimates do --- and don't show}
\small
The composed score is above the tuned guard on transfer data for every model in this fixed panel.
It does \textbf{not} beat every untuned base: it is slightly below SmolLM3's base
($0.907$ vs $0.910$) and below the strongest model, Qwen3-4B ($0.914$ vs $0.944$).
So composing can \textbf{recover much of the damage fine-tuning did} without dominating every
base or SFT result (\Cref{tab:composeperm}).
\end{takeaway}

Read \Cref{tab:composeperm} across each row: the composed column is always above the tuned one.
It exceeds the base for Qwen2.5 and SmolLM2 but remains just below the base for SmolLM3 and
Qwen3-4B. The per-model effects are mixed. Because same-row base AP is also part of the
composed-minus-base contrast, this table cannot by itself show that base strength predicts recovery.

\begin{table}[t]
\centering\small
\caption{Per-model new-data (transfer) macro-AP. Composition exceeds SFT for all four models,
but its comparison with base is mixed; it is therefore not a universal or Pareto improvement.}
\label{tab:composeperm}
\begin{tabular}{lrrrr}
\toprule
Model (new-data strength of base) & Untuned base & Tuned & Composed & Comp. $-$ base \\
\midrule
SmolLM2-1.7B (weakest) & 0.790 & 0.830 & 0.857 & $+0.066$ \\
Qwen2.5-1.5B & 0.819 & 0.780 & 0.855 & $+0.036$ \\
SmolLM3-3B & 0.910 & 0.823 & 0.907 & $-0.003$ \\
Qwen3-4B (strongest) & 0.944 & 0.794 & 0.914 & $-0.029$ \\
\bottomrule
\end{tabular}
\end{table}

\begin{background}{is this real, or just a trick of averaging?}
\small
A fair worry is whether the result reflects a generic averaging effect. Two single-permutation
diagnostics probe, but do not settle, that question. (1) Scrambling the tuned guard's scores so
they no longer track safe-vs-unsafe makes the observed boost negative. (2) Shuffling tuned scores
within each gold class preserves their marginal signal while breaking row-level pairing with the
base; the boost survives. These diagnostics narrow explanations under those perturbations, but
they do not distinguish retaining the base from generic two-pass ensembling, establish a
mechanism, or provide prospective evidence. A same-inference-cost SFT+SFT control is still required.
\end{background}

\paragraph{Limits of this Paper B preview.} Reusing Paper A's clean-v2 retrospective scored rows makes
this exploratory, not an independent replication or prospective test. Composing improves the
observed \emph{ranking}, not necessarily \emph{calibration}; transfer false-alarm behavior needs
separately held-out threshold evaluation. It also costs two passes per prompt. Independent,
prospectively collected evaluation data, a same-inference-cost SFT+SFT control, and a competence
measure taken on data separate from the recovery outcome are needed before treating composition
as a general remedy or decision rule.

% ============================================================
\section{What this means in practice}
% ============================================================
\begin{takeaway}{Practitioner takeaways}
\small
\begin{itemize}
\item \textbf{Fine-tune when your live traffic resembles your training sources.}
On these represented datasets, the observed gain is large (recall $13.0\%\!\to\!76.9\%$); validate
the effect on your own traffic.
\item \textbf{Check for erosion on dataset-held-out attacks.} On the transfer benchmarks,
fine-tuning shows more false alarms and an observed drop on hard harmful examples.
\item \textbf{Always test on datasets the model never trained on.} A single trained-on
benchmark will crown fine-tuning the winner and hide the transfer cost. Report \emph{both}
groups.
\item \textbf{Do not route by starting strength from this panel.} The four checkpoint effects are
heterogeneous, but same-row base-versus-change comparisons are coupled. Test every base/adapted
pair on target transfer data; validate any competence rule on an independent development measure
and a prospectively locked outcome.
\item \textbf{Treat composition as a Paper B hypothesis.} Keeping the untuned base in the
decision recovers some transfer ranking in the preliminary fixed-panel analysis
(\Cref{sec:fix}), but it needs independent prospective evaluation.
\end{itemize}
\end{takeaway}

% ============================================================
\section{Limitations (what not to over-read)}
% ============================================================
\begin{itemize}
\item \textbf{Retrospective fixed-panel evidence.} The analysis is \emph{estimation-only}.
Re-running a locked pipeline can verify execution but cannot make previously inspected transfer
sets prospective. Confirmatory claims require a new sealed cohort and a prospectively locked plan.
\item \textbf{Balanced test pools overstate real-world precision.} Test sets are ${\sim}50/50$
safe/unsafe; real traffic has far fewer unsafe prompts, so precision/AP look rosier here than
in production.
\item \textbf{Transfer sets were not truly sealed.} They were inspected during development, so
``transfer'' means dataset-held-out, not never-before-seen.
\item \textbf{Only four models, two families (Qwen, SmolLM), 1.5--4B.} Conclusions are about
\emph{this panel}; the base-competence interpretation rests on four mathematically coupled,
same-row points and remains a hypothesis.
\item \textbf{Mixed native policies.} The four transfer datasets encode different definitions
of ``unsafe'' (harm vs.\ refusal vs.\ jailbreak), so their macro-average blends distinct
policies.
\item \textbf{Prompt-only, English.} No response, tool-call, or multi-turn moderation.
\item \textbf{The composition analysis (\Cref{sec:fix}) belongs to preliminary Paper B work.}
It reuses clean-v2 retrospective Paper A scores, does not establish transfer calibration, and costs two
passes per prompt. It requires an independent prospective evaluation.
\end{itemize}

% ============================================================
\section{Conclusion}
% ============================================================
In this fixed panel, fine-tuning \emph{specializes} a small safety guard: the analysis estimates
large gains on data resembling its training sources (represented macro-AP $+0.32$, everyone to
${\approx}0.98$), a small average loss on dataset-held-out transfer benchmarks ($-0.06$), and
an observed HarmBench stress-set drop (recall $-18.0$ points). The average transfer number hides
checkpoint heterogeneity --- one gain, one interval spanning zero, and two losses --- and 15 of 20 fine-tunes land in the
``better-on-familiar, worse-on-new'' quadrant. So the benchmark you choose to report decides
whether fine-tuning looks like a win. Report \emph{both} groups with conditional fixed-panel
resampling intervals and treat these clean-v2 numbers as retrospective estimates, not prospective
or confirmatory effects. A separate preliminary
Paper B analysis (\Cref{sec:fix}) suggests composition as a mitigation candidate; it is not a
Paper A conclusion and still needs an independent prospective test.

% ============================================================
\appendix
\section{Per-seed changes}
% ============================================================
Table~\ref{tab:seeds} lists the paired change for every (checkpoint, seed), the raw material
behind the panel aggregates and Figure~\ref{fig:plane}. This spread describes run-to-run
optimization variation with fixed data order under one recipe; it does not measure uncertainty
over new checkpoints, datasets, or recipes.

\begin{table}[h]
\centering\small
\caption{Represented and transfer AP change ($\Delta$) for each of the 20 fine-tunes.}
\label{tab:seeds}
\begin{tabular}{l rrrrr}
\toprule
& \multicolumn{5}{c}{seed} \\
\cmidrule(lr){2-6}
Checkpoint (group) & 42 & 43 & 44 & 45 & 46 \\
\midrule
Qwen2.5-1.5B (rep.)  & $+0.3569$ & $+0.3535$ & $+0.3506$ & $+0.3558$ & $+0.3551$ \\
Qwen2.5-1.5B (tr.)   & $+0.0082$ & $-0.0261$ & $-0.0545$ & $-0.0438$ & $-0.0786$ \\
SmolLM2-1.7B (rep.)  & $+0.5304$ & $+0.5273$ & $+0.5253$ & $+0.5309$ & $+0.5272$ \\
SmolLM2-1.7B (tr.)   & $+0.0805$ & $+0.0720$ & $+0.0250$ & $+0.0238$ & $-0.0013$ \\
SmolLM3-3B (rep.)   & $+0.3045$ & $+0.3253$ & $+0.3051$ & $+0.3146$ & $+0.3156$ \\
SmolLM3-3B (tr.)    & $-0.0812$ & $-0.0925$ & $-0.0629$ & $-0.0923$ & $-0.1055$ \\
Qwen3-4B (rep.)    & $+0.0953$ & $+0.0936$ & $+0.1038$ & $+0.0997$ & $+0.0981$ \\
Qwen3-4B (tr.)     & $-0.1222$ & $-0.1536$ & $-0.0949$ & $-0.2291$ & $-0.1497$ \\
\bottomrule
\end{tabular}
\end{table}

\section{Provenance and reproducibility}
The scored data table stores per-row \emph{content hashes and model logits only --- never raw
prompt text} --- so third-party benchmark text is not redistributed, yet every metric is
recomputable. Metrics come from one shared, tie-aware module (no ad-hoc AP loop). Source
datasets are pinned by revision; near-duplicate ``families'' are detected with a fixed MinHash
so the split cannot change with the environment. Full identifiers, hashes, licenses, and
seed-level outputs are in \code{artifacts/paper\_a\_sft\_v2/}.

\section*{Related work}
This edition draws on: purpose-built prompt guards --- Llama Guard~\cite{inan2023llamaguard},
ShieldGemma~\cite{zeng2024shieldgemma}, Granite Guardian~\cite{padhi2024graniteguardian},
WildGuard~\cite{han2024wildguard}, and the parameter-efficient LoRA-Guard~\cite{elesedy2024loraguard};
evidence that fine-tuning can erode guardrails and that guards lean on surface heuristics
\cite{hsiung2025guardrails,li2026surfaceheuristics,bassani2025robustness}, which our
specialization measurement quantifies per checkpoint. The preliminary Paper B composition analysis (\Cref{sec:fix})
relates to \emph{weight-space} robust fine-tuning and model soups~\cite{wortsman2022wiseft,wortsman2022modelsoups}
and to calibrated ensembles under distribution shift~\cite{kumar2022calibrated}; we combine
\emph{scores} (output-space), needing no shared weights. Other tuning objectives (DPO, GRPO) are
a separate lever we do not vary here~\cite{rafailov2023dpo,shao2024deepseekmath,vennemeyer2026objective}.
Guard calibration and evaluation context: \cite{guo2017calibration,liu2025calibration,bassani2024guardbench}.
For the full formal protocol and the lock/audit contract, see the formal paper under \code{papers/finetuning-specialization/}.

\bibliographystyle{plain}
\bibliography{refs}

\end{document}
